\documentclass[11pt]{article}
\usepackage[margin=1in]{geometry}
\usepackage{amsmath,amssymb,amsthm,mathtools}
\usepackage{enumitem}
\usepackage{booktabs}
\usepackage{hyperref}
\usepackage{xcolor}

\newtheorem{theorem}{Theorem}
\newtheorem{proposition}{Proposition}
\newtheorem{lemma}{Lemma}
\newtheorem{corollary}{Corollary}
\theoremstyle{definition}
\newtheorem{remark}{Remark}
\newtheorem{definition}{Definition}

% Notation
\newcommand{\CP}{\mathcal{C}^{*}}
\newcommand{\COP}{\mathcal{C}}
\newcommand{\R}{\mathbb{R}}
\newcommand{\N}{\mathcal{N}}
\newcommand{\inner}[2]{\langle #1, #2 \rangle}
\newcommand{\Tr}{\operatorname{Tr}}
\newcommand{\diag}{\operatorname{diag}}
\newcommand{\conv}{\operatorname{conv}}
\newcommand{\xcom}{x^{\mathrm{cpl}}}
\newcommand{\xrest}{x^{\mathrm{rest}}}
\newcommand{\Acom}{A^{\mathrm{cpl}}}
\newcommand{\feas}{\mathcal{X}}
\newcommand{\feaseq}{\mathcal{X}^{\mathrm{eq}}}
\newcommand{\feasbar}{\bar{\mathcal{X}}^{\mathrm{eq}}}
\newcommand{\liftset}{\mathcal{F}}
\newcommand{\vmip}{v^{\mathrm{MIP}}}
\newcommand{\vchp}{v^{\mathrm{CHP}}}
\newcommand{\vlr}{v^{\mathrm{LR}}}
\newcommand{\Tech}{\mathcal{T}}

\title{Completely Positive Reformulation of the\\
Mixed-Integer Energy Community Game:\\
A Technology-Modular Construction}
\author{Working note}
\date{\today}

\begin{document}
\maketitle

\begin{abstract}
We construct the completely positive program~(CPP) reformulation of the
mixed-integer community game from the companion preprint.  The reformulation
is organised by \emph{technology module}---renewable generation, energy
storage, electrolyzer with piecewise-linear efficiency, heat pump with unit
commitment, and demand---so that any player with any combination of
technologies receives its CPP lift by composing the relevant modules.  The
only inter-player coupling retained here is the community balance (one
homogeneous equality per carrier per period); reserve adequacy and peak
penalty couplings are deferred to future work.  We verify Burer's exactness
conditions module by module, write the Lagrangian dual (COP), confirm that
cross-block vanishing reduces the relaxed game to a classical linear
production game, and restate the main core-existence theorem.  Convention:
profit maximisation; all community balance signs unified to
$\sum_i(i^{\mathrm{com}}_{k,t}-e^{\mathrm{com}}_{k,t})=0$.
\end{abstract}

%======================================================================
\section{The general MIP community game}
\label{sec:mip}
%======================================================================

\subsection{Players and technology modules}

Let $\N=\{1,\dots,n\}$ be the prosumers.  Each player $i$ owns a
technology set
\[
\Tech_i\subseteq
\{\textsc{res},\textsc{ess}_E,\textsc{ess}_G,\textsc{ess}_H,
  \textsc{els},\textsc{hp},
  \textsc{d}_E,\textsc{d}_G,\textsc{d}_H,
  \textsc{fl}_E,\textsc{fl}_G,\textsc{fl}_H\},
\]
where the symbols denote:
\begin{center}\small
\begin{tabular}{@{}ll@{}}
\toprule
Module & Description\\
\midrule
\textsc{res} & Renewable (e.g.\ wind) generation \\
$\textsc{ess}_k$ & Energy storage, carrier $k\in\{E,G,H\}$\\
\textsc{els} & Electrolyzer with piecewise-linear efficiency \\
\textsc{hp} & Heat pump with unit commitment \\
$\textsc{d}_k$ & Non-flexible demand, carrier $k$\\
$\textsc{fl}_k$ & Flexible demand, carrier $k$ (e.g.\ electrolyzer/HP elec.\ input)\\
\bottomrule
\end{tabular}
\end{center}
A player's decision variables $x_i=(\xcom_i,\xrest_i)$ and
per-player feasible region $\feas_i$ are assembled by taking the union
of variables and the intersection of constraint sets across all modules
in~$\Tech_i$, together with the per-carrier energy balance.

\subsection{Variable classification per technology}
\label{sec:vars}

We partition each player's variables into \emph{coupling variables}
$\xcom_i$ (those appearing in the inter-player community balance) and
\emph{rest variables} $\xrest_i$ (per-player only, including binaries).
Table~\ref{tab:tech-vars} catalogues the variables contributed by each
module, per time period~$t$.

\begin{table}[h]
\centering\small
\begin{tabular}{@{}l l l l@{}}
\toprule
Module & Coupling (per $t$) & Rest continuous (per $t$) & Rest binary (per $t$)\\
\midrule
\textsc{res}
  & $e^{\mathrm{com},E}_{i,t}$
  & $e^{\mathrm{gri},E}_{i,t},\; p^{\mathrm{res}}_{i,t}$
  & ---\\[3pt]
\textsc{els}
  & $i^{\mathrm{com},E}_{i,t},\; i^{\mathrm{com},G}_{i,t},\; e^{\mathrm{com},G}_{i,t}$
  & \makecell[tl]{$i^{\mathrm{gri},E}_{i,t},\; e^{\mathrm{gri},G}_{i,t},\; p^{\mathrm{els}}_{i,t}$,\\
                  $d^{\mathrm{els}}_{i,t,s}\ (s=1,\dots,N_s)$,\\
                  $d^{\mathrm{fl},E}_{i,t}$}
  & \makecell[tl]{$z^{\mathrm{on}}_{i,t},\; z^{\mathrm{off}}_{i,t},\; z^{\mathrm{sb}}_{i,t}$,\\
                  $z^{\mathrm{su}}_{i,t},\; z^{\mathrm{sd}}_{i,t}$,\\
                  $z^{\mathrm{seg}}_{i,t,s}\ (s=1,\dots,N_s)$}\\[3pt]
\textsc{hp}
  & $i^{\mathrm{com},E}_{i,t},\; i^{\mathrm{com},H}_{i,t},\; e^{\mathrm{com},H}_{i,t}$
  & \makecell[tl]{$i^{\mathrm{gri},E}_{i,t},\; e^{\mathrm{gri},H}_{i,t},\; p^{\mathrm{hp}}_{i,t}$,\\
                  $d^{\mathrm{fl},E}_{i,t}$}
  & $z^{\mathrm{on}}_{i,t},\; z^{\mathrm{su}}_{i,t},\; z^{\mathrm{sd}}_{i,t}$\\[3pt]
$\textsc{ess}_k$
  & ---
  & $b^{\mathrm{ch},k}_{i,t},\; b^{\mathrm{dis},k}_{i,t},\; s^{k}_{i,t}$
  & ---\\[3pt]
$\textsc{d}_k$
  & $i^{\mathrm{com},k}_{i,t}$
  & $i^{\mathrm{gri},k}_{i,t},\; d^{\mathrm{nfl},k}_{i,t}$
  & ---\\[3pt]
$\textsc{fl}_k$
  & $i^{\mathrm{com},k}_{i,t}$
  & $i^{\mathrm{gri},k}_{i,t},\; d^{\mathrm{fl},k}_{i,t}$
  & ---\\
\bottomrule
\end{tabular}
\caption{Variables per technology module and period.  A player $i$ with
$\Tech_i=\{\textsc{els},\textsc{ess}_G,\textsc{fl}_E\}$ takes the union
of the corresponding rows; coupling variables shared between modules
(e.g.\ $i^{\mathrm{com},E}$ from both \textsc{els} and \textsc{fl}$_E$)
appear once.}
\label{tab:tech-vars}
\end{table}

\subsection{Coupling constraints (homogeneous community balance)}

For each carrier $k\in\{E,G,H\}$ and period $t\in\{0,\dots,T-1\}$:
\begin{equation}
\sum_{i\in S}\bigl(i^{\mathrm{com},k}_{i,t}-e^{\mathrm{com},k}_{i,t}\bigr)=0.
\label{eq:community-balance}
\end{equation}
The right-hand side is \emph{zero} for every row (conservation law:
internal community trade nets to zero).  Written compactly,
\begin{equation}
\sum_{i\in S}\Acom_i\,\xcom_i=0,
\qquad R=|K|\cdot T\text{ rows, all RHS }=0.
\label{eq:coupling-compact}
\end{equation}
The matrix $\Acom_i\in\R^{R\times n_i^{\mathrm{cpl}}}$ has entry $+1$
on $i^{\mathrm{com},k}_{i,t}$ and $-1$ on $e^{\mathrm{com},k}_{i,t}$
in the row for carrier $k$, period $t$; zero elsewhere.

\begin{remark}[Sign convention]
\label{rem:sign}
In the code (\texttt{compact\_utility.py}), the hydrogen balance is
written $\sum_i(e^{\mathrm{com},G}-i^{\mathrm{com},G})=0$ (reversed
signs).  We unify all carriers to the
$\sum_i(i^{\mathrm{com},k}-e^{\mathrm{com},k})=0$ convention; this
amounts to redefining $\Acom_i$ for the hydrogen rows by swapping
signs, with no effect on the feasible set.
\end{remark}

\subsection{Per-player constraints (by technology module)}

The feasible region of player $i$ is
\begin{equation}
\feas_i=\bigcap_{\tau\in\Tech_i}\feas_i^{\tau}
  \cap\bigl\{\text{energy balance equalities}\bigr\},
\label{eq:feasi}
\end{equation}
where each $\feas_i^{\tau}$ collects the constraints of module $\tau$
(detailed in Section~\ref{sec:catalog}).  All constraints in $\feas_i$
involve \emph{only} player $i$'s variables.  $\feas_i$ is mixed-integer
and bounded (all variables have finite upper bounds from capacities,
demands, or storage sizes).

\subsection{The coalition game}

\begin{equation}
\vmip(S)=\max\Bigl\{\sum_{i\in S}w_i^{\top}x_i\;:\;
\sum_{i\in S}\Acom_i\xcom_i=0,\;x_i\in\feas_i\ \forall i\in S\Bigr\},
\label{eq:vmip}
\end{equation}
where $w_i$ is the profit coefficient vector for player $i$ (revenue
minus cost in the objective; note the code minimises cost, so
$w_i=-c_i$).

%======================================================================
\section{Technology constraint catalog}
\label{sec:catalog}
%======================================================================

For each module we give: (i)~the original constraints as implemented in
\texttt{compact\_utility.py}; (ii)~the equality form with nonnegative
slacks; (iii)~the resulting rows of $A_ix_i=b_i$; (iv)~the RLT
identities (diagonal of the squared system); (v)~binary diagonal
conditions.  The generic conversion rules are:
\begin{itemize}[leftmargin=1.2em]
\item Inequality $g(x)\le 0$: add slack $\sigma\ge0$, write
  $g(x)+\sigma=0$.
\item Upper bound $x\le\bar x$: add slack $\sigma\ge0$, write
  $x+\sigma=\bar x$.
\item Binary $z\in\{0,1\}$: add complement $\sigma\ge0$, write
  $z+\sigma=1$; the quadratic identity $z-z^2=0$ is handled at the
  lifting step (Step~2).
\end{itemize}

\subsection{Energy balance (per player, per carrier, per period)}
\label{sec:balance}

Each player has one balance per carrier it participates in.  For carrier
$k$ and period~$t$:
\begin{equation}
\underbrace{i^{\mathrm{gri},k}_{i,t}-e^{\mathrm{gri},k}_{i,t}
+i^{\mathrm{com},k}_{i,t}-e^{\mathrm{com},k}_{i,t}}_{\text{trading}}
+\underbrace{p^{k}_{i,t}+b^{\mathrm{dis},k}_{i,t}
  -b^{\mathrm{ch},k}_{i,t}}_{\text{production/storage}}
=\underbrace{d^{\mathrm{nfl},k}_{i,t}+d^{\mathrm{fl},k}_{i,t}}_{\text{demand}},
\label{eq:energy-balance}
\end{equation}
where terms absent for a given player (because $\Tech_i$ lacks the
module) are zero.  This is already an equality; it contributes one row
of $A_ix_i=b_i$ per carrier-period pair in which player $i$ participates.

\subsection{Renewable generation (\textsc{res})}
\label{sec:res}

\paragraph{Original constraints.}
\begin{equation}
0\le p^{\mathrm{res}}_{i,t}\le\overline{P}^{\mathrm{res}}_{i,t},
\label{eq:res-orig}
\end{equation}
where $\overline{P}^{\mathrm{res}}_{i,t}$ is the time-dependent
renewable availability (wind capacity factor $\times$ installed
capacity).

\paragraph{Equality form.}  Add slack $\sigma^{\mathrm{res}}_{i,t}\ge0$:
\begin{equation}
p^{\mathrm{res}}_{i,t}+\sigma^{\mathrm{res}}_{i,t}
  =\overline{P}^{\mathrm{res}}_{i,t}.
\label{eq:res-eq}
\end{equation}

\paragraph{RLT identity (squaring \eqref{eq:res-eq}).}
\begin{equation}
(X_{ii})_{p^{\mathrm{res}},p^{\mathrm{res}}}
+2(X_{ii})_{p^{\mathrm{res}},\sigma^{\mathrm{res}}}
+(X_{ii})_{\sigma^{\mathrm{res}},\sigma^{\mathrm{res}}}
=\bigl(\overline{P}^{\mathrm{res}}_{i,t}\bigr)^{2}.
\label{eq:res-rlt}
\end{equation}

\paragraph{Binaries.} None.

\paragraph{Objective.}
$w_i$ entries: $-c^{\mathrm{res}}_i$ on $p^{\mathrm{res}}_{i,t}$
(production cost, sign-flipped to profit); grid export revenue
$+\pi^{E,\mathrm{exp}}_t$ on $e^{\mathrm{gri},E}_{i,t}$; grid export
capacity bound $e^{\mathrm{gri},E}_{i,t}\le\overline{E}^{E}_i$ (slack
$\sigma^{E,\mathrm{exp}}_{i,t}$, row
$e^{\mathrm{gri},E}+\sigma^{E,\mathrm{exp}}=\overline{E}^{E}_i$).

\subsection{Energy storage (\texorpdfstring{$\textsc{ess}_k$}{ESS\_k})}
\label{sec:ess}

\paragraph{Original constraints (per period $t$, carrier $k$).}
\begin{align}
s^{k}_{i,t}&=s^{k}_{i,t-1}+\eta^{\mathrm{ch}}_k\,b^{\mathrm{ch},k}_{i,t}
  -\tfrac{1}{\eta^{\mathrm{dis}}_k}\,b^{\mathrm{dis},k}_{i,t},
  &\text{(SOC transition)}\label{eq:soc}\\
s^{k}_{i,t_0}&=\mathrm{SOC}^{\mathrm{init}}_k,
  &\text{(initial SOC)}\label{eq:soc-init}\\
s^{k}_{i,0}&=s^{k}_{i,T-1}+\eta^{\mathrm{ch}}_k\,b^{\mathrm{ch},k}_{i,0}
  -\tfrac{1}{\eta^{\mathrm{dis}}_k}\,b^{\mathrm{dis},k}_{i,0},
  &\text{(wrap-around)}\label{eq:soc-wrap}\\
0\le s^{k}_{i,t}&\le\overline{S}^{k}_i,\quad
0\le b^{\mathrm{ch},k}_{i,t}\le\overline{P}^{k}_i,\quad
0\le b^{\mathrm{dis},k}_{i,t}\le\overline{P}^{k}_i.
  &\text{(bounds)}\label{eq:ess-bounds}
\end{align}

\paragraph{Equality form.}
The SOC transitions \eqref{eq:soc}--\eqref{eq:soc-wrap} and the initial
condition~\eqref{eq:soc-init} are already equalities.  Each upper bound
in~\eqref{eq:ess-bounds} adds a slack:
\begin{align}
s^{k}_{i,t}+\sigma^{s,k}_{i,t}&=\overline{S}^{k}_i,\label{eq:ess-eq-s}\\
b^{\mathrm{ch},k}_{i,t}+\sigma^{\mathrm{ch},k}_{i,t}&=\overline{P}^{k}_i,
\label{eq:ess-eq-ch}\\
b^{\mathrm{dis},k}_{i,t}+\sigma^{\mathrm{dis},k}_{i,t}&=\overline{P}^{k}_i.
\label{eq:ess-eq-dis}
\end{align}

\paragraph{RLT identities.}
One per equality row.  For example, squaring
\eqref{eq:ess-eq-ch}:
\begin{equation}
(X_{ii})_{b^{\mathrm{ch}},b^{\mathrm{ch}}}
+2(X_{ii})_{b^{\mathrm{ch}},\sigma^{\mathrm{ch}}}
+(X_{ii})_{\sigma^{\mathrm{ch}},\sigma^{\mathrm{ch}}}
=(\overline{P}^{k}_i)^{2}.
\label{eq:ess-rlt-ch}
\end{equation}
The SOC-transition rows produce multivariate RLT identities of the form
$\inner{a_ja_j^{\top}}{X_{ii}}=b_j^{2}$ where $a_j$ is the row of
$A_i$ corresponding to~\eqref{eq:soc}.

\paragraph{Binaries.} None.

\paragraph{Objective.}
$w_i$ entries: $-c^{\mathrm{sto},k}_i\eta^{\mathrm{ch}}_k$ on
$b^{\mathrm{ch},k}_{i,t}$;
$-c^{\mathrm{sto},k}_i/\eta^{\mathrm{dis}}_k$ on
$b^{\mathrm{dis},k}_{i,t}$ (storage operation costs, sign-flipped).

\subsection{Electrolyzer with piecewise-linear efficiency (\textsc{els})}
\label{sec:els}

This is the most complex module: it carries $5+N_s$ binary variables per
period (for $N_s$ PWL segments, typically $N_s=6$).

\paragraph{Original constraints (per period $t$).}

\emph{State selection} (exactly one of on/off/standby):
\begin{equation}
z^{\mathrm{on}}_{i,t}+z^{\mathrm{off}}_{i,t}+z^{\mathrm{sb}}_{i,t}=1.
\label{eq:els-state}
\end{equation}

\emph{PWL production curve} ($N_s$ segments with slope $a_s$ and
intercept $b_s$):
\begin{equation}
p^{\mathrm{els}}_{i,t}=\sum_{s=1}^{N_s}\bigl[
  a_s\,d^{\mathrm{els}}_{i,t,s}+b_s\,z^{\mathrm{seg}}_{i,t,s}\bigr].
\label{eq:els-production}
\end{equation}

\emph{Segment selection}:
\begin{equation}
z^{\mathrm{on}}_{i,t}=\sum_{s=1}^{N_s}z^{\mathrm{seg}}_{i,t,s}.
\label{eq:els-seg-sel}
\end{equation}

\emph{Segment bounds} (for each $s=1,\dots,N_s$):
\begin{equation}
\underline{p}_s\,z^{\mathrm{seg}}_{i,t,s}
  \le d^{\mathrm{els}}_{i,t,s}
  \le\overline{p}_s\,z^{\mathrm{seg}}_{i,t,s},
\label{eq:els-seg-bounds}
\end{equation}
where $\underline{p}_s,\overline{p}_s$ are the breakpoints of segment $s$.

\emph{Power consumption coupling}:
\begin{equation}
d^{\mathrm{fl},E}_{i,t}=\sum_{s=1}^{N_s}d^{\mathrm{els}}_{i,t,s}
  +c^{\mathrm{sb}}\,\overline{C}^{\mathrm{els}}_i\,z^{\mathrm{sb}}_{i,t}.
\label{eq:els-flex}
\end{equation}

\emph{Total power bounds}:
\begin{align}
d^{\mathrm{fl},E}_{i,t}&\ge c^{\min}\,\overline{C}^{\mathrm{els}}_i\,
  z^{\mathrm{on}}_{i,t}
  +c^{\mathrm{sb}}\,\overline{C}^{\mathrm{els}}_i\,z^{\mathrm{sb}}_{i,t},
\label{eq:els-min-total}\\
d^{\mathrm{fl},E}_{i,t}&\le c^{\max}\,\overline{C}^{\mathrm{els}}_i\,
  z^{\mathrm{on}}_{i,t}
  +c^{\mathrm{sb}}\,\overline{C}^{\mathrm{els}}_i\,z^{\mathrm{sb}}_{i,t}.
\label{eq:els-max-total}
\end{align}

\emph{Startup/shutdown logic} (for $t\neq t_0$; write $t^{-}$ for the
preceding period, with wrap-around $0^{-}=T{-}1$):
\begin{align}
z^{\mathrm{su}}_{i,t}&\ge z^{\mathrm{on}}_{i,t}
  +z^{\mathrm{off}}_{i,t^-}+z^{\mathrm{sb}}_{i,t}-1,
\label{eq:els-su}\\
z^{\mathrm{sd}}_{i,t}&\ge z^{\mathrm{on}}_{i,t^-}
  +z^{\mathrm{off}}_{i,t}+z^{\mathrm{sb}}_{i,t^-}-1,
\label{eq:els-sd}\\
z^{\mathrm{su}}_{i,t}+z^{\mathrm{sd}}_{i,t}&\le1,
\label{eq:els-conflict}\\
z^{\mathrm{off}}_{i,t^-}+z^{\mathrm{sb}}_{i,t}&\le1.
\label{eq:els-off-sb}
\end{align}

\emph{Minimum downtime} ($\Delta T^{\mathrm{down}}$ periods):
\begin{equation}
\sum_{\tau=t}^{t+\Delta T^{\mathrm{down}}-1}z^{\mathrm{off}}_{i,\tau}
  \ge\Delta T^{\mathrm{down}}\,z^{\mathrm{sd}}_{i,t}.
\label{eq:els-mindown}
\end{equation}

\emph{Initial state}: $z^{\mathrm{off}}_{i,t_0}=1$.

\paragraph{Equality form.}

Equalities~\eqref{eq:els-state},~\eqref{eq:els-production},~\eqref{eq:els-seg-sel},~\eqref{eq:els-flex}
remain as-is.  Each inequality adds a nonneg slack; we illustrate the
pattern.  For the segment lower bound~\eqref{eq:els-seg-bounds}:
\begin{equation}
d^{\mathrm{els}}_{i,t,s}-\underline{p}_s\,z^{\mathrm{seg}}_{i,t,s}
  -\sigma^{\mathrm{lb}}_{i,t,s}=0,
\qquad\sigma^{\mathrm{lb}}_{i,t,s}\ge0.
\label{eq:els-seg-lb-eq}
\end{equation}
For the startup~\eqref{eq:els-su}:
\begin{equation}
z^{\mathrm{su}}_{i,t}-z^{\mathrm{on}}_{i,t}
-z^{\mathrm{off}}_{i,t^-}-z^{\mathrm{sb}}_{i,t}+1
-\sigma^{\mathrm{su}}_{i,t}=0,
\qquad\sigma^{\mathrm{su}}_{i,t}\ge0.
\label{eq:els-su-eq}
\end{equation}
And so on for each inequality row.

Each binary $z\in\{z^{\mathrm{on}},z^{\mathrm{off}},z^{\mathrm{sb}},
z^{\mathrm{su}},z^{\mathrm{sd}},z^{\mathrm{seg}}_1,\dots,
z^{\mathrm{seg}}_{N_s}\}$ produces a complement slack:
\begin{equation}
z+\sigma^z=1,\qquad z\ge0,\;\sigma^z\ge0.
\label{eq:els-bin-complement}
\end{equation}
The quadratic identity $z-z^2=0$ is handled in the lifting
(Section~\ref{sec:lifting}).

\paragraph{Count per period.}
$5+N_s$ binary complement rows, plus $2N_s$ segment bound rows, plus 4
startup/shutdown rows, plus 1 conflict, plus 1 off-to-standby, plus
$\sim\!\Delta T^{\mathrm{down}}$ downtime rows, plus 2 total-power
bounds.  The exact count depends on the horizon and wrap-around handling
but is $O(N_s+\Delta T^{\mathrm{down}})$ per period.

\paragraph{RLT identities (illustrative examples).}

Squaring the state selection~\eqref{eq:els-state}:
\begin{equation}
(X_{ii})_{z^{\mathrm{on}},z^{\mathrm{on}}}
+(X_{ii})_{z^{\mathrm{off}},z^{\mathrm{off}}}
+(X_{ii})_{z^{\mathrm{sb}},z^{\mathrm{sb}}}
+2(X_{ii})_{z^{\mathrm{on}},z^{\mathrm{off}}}
+2(X_{ii})_{z^{\mathrm{on}},z^{\mathrm{sb}}}
+2(X_{ii})_{z^{\mathrm{off}},z^{\mathrm{sb}}}
=1.
\label{eq:els-rlt-state}
\end{equation}
Squaring a binary complement $z^{\mathrm{on}}+\sigma^{\mathrm{on}}=1$:
\begin{equation}
(X_{ii})_{z^{\mathrm{on}},z^{\mathrm{on}}}
+2(X_{ii})_{z^{\mathrm{on}},\sigma^{\mathrm{on}}}
+(X_{ii})_{\sigma^{\mathrm{on}},\sigma^{\mathrm{on}}}=1.
\label{eq:els-rlt-comp}
\end{equation}
Combining~\eqref{eq:els-rlt-state} with the binary diagonal
$(X_{ii})_{z^{\mathrm{on}},z^{\mathrm{on}}}=z^{\mathrm{on}}$ (see
below) gives cross-moment bounds: e.g.,
$(X_{ii})_{z^{\mathrm{on}},\sigma^{\mathrm{on}}}
=\tfrac{1}{2}(1-z^{\mathrm{on}}-(X_{ii})_{\sigma^{\mathrm{on}},\sigma^{\mathrm{on}}})$.

\paragraph{Binary diagonal (per period $t$).}
For each $m\in B_i$:
\begin{equation}
(X_{ii})_{m,m}=z_{i,m}.
\label{eq:els-bindiag}
\end{equation}
This gives $5+N_s$ conditions per period.  For $N_s=6$: $11$ conditions
$\times T=264$ total (in the standard $T=24$ model).

\subsection{Heat pump with unit commitment (\textsc{hp})}
\label{sec:hp}

\paragraph{Original constraints (per period $t$).}

\emph{COP coupling}:
\begin{equation}
\eta^{\mathrm{COP}}_i\,d^{\mathrm{fl},E}_{i,t}=p^{\mathrm{hp}}_{i,t}.
\label{eq:hp-cop}
\end{equation}

\emph{State transition} ($t\neq0$; $t^{-}$ as before):
\begin{equation}
z^{\mathrm{su}}_{i,t}-z^{\mathrm{sd}}_{i,t}
  -z^{\mathrm{on}}_{i,t}+z^{\mathrm{on}}_{i,t^{-}}=0.
\label{eq:hp-state}
\end{equation}

\emph{Mutual exclusion}:
\begin{equation}
z^{\mathrm{su}}_{i,t}+z^{\mathrm{sd}}_{i,t}\le1.
\label{eq:hp-conflict}
\end{equation}

\emph{Power bounds}:
\begin{align}
p^{\mathrm{hp}}_{i,t}&\ge c^{\min}_{\mathrm{hp}}\,
  \overline{C}^{\mathrm{hp}}_i\,z^{\mathrm{on}}_{i,t},
\label{eq:hp-min}\\
p^{\mathrm{hp}}_{i,t}&\le c^{\max}_{\mathrm{hp}}\,
  \overline{C}^{\mathrm{hp}}_i\,z^{\mathrm{on}}_{i,t}.
\label{eq:hp-max}
\end{align}

\emph{Ramp rates} ($t\neq0$):
\begin{align}
p^{\mathrm{hp}}_{i,t}-p^{\mathrm{hp}}_{i,t^{-}}
  &\le R^{\mathrm{up}}\,\overline{C}^{\mathrm{hp}}_i\,z^{\mathrm{on}}_{i,t}
  +R^{\mathrm{su}}\,\overline{C}^{\mathrm{hp}}_i\,z^{\mathrm{su}}_{i,t},
\label{eq:hp-ramp-up}\\
p^{\mathrm{hp}}_{i,t^{-}}-p^{\mathrm{hp}}_{i,t}
  &\le R^{\mathrm{down}}\,\overline{C}^{\mathrm{hp}}_i\,z^{\mathrm{on}}_{i,t}
  +R^{\mathrm{sd}}\,\overline{C}^{\mathrm{hp}}_i\,z^{\mathrm{sd}}_{i,t}.
\label{eq:hp-ramp-down}
\end{align}

\emph{Initial state}: $z^{\mathrm{on}}_{i,t_0}=0$.

\paragraph{Equality form.}
The COP coupling~\eqref{eq:hp-cop} and state
transition~\eqref{eq:hp-state} are already equalities.  Each of
\eqref{eq:hp-conflict}--\eqref{eq:hp-ramp-down} adds a nonneg slack.
Each binary
$z\in\{z^{\mathrm{on}},z^{\mathrm{su}},z^{\mathrm{sd}}\}$ produces a
complement slack $z+\sigma^z=1$.

\paragraph{Binary diagonal (per period).}
3 conditions: $(X_{ii})_{z^{\mathrm{on}},z^{\mathrm{on}}}=z^{\mathrm{on}}$,
$(X_{ii})_{z^{\mathrm{su}},z^{\mathrm{su}}}=z^{\mathrm{su}}$,
$(X_{ii})_{z^{\mathrm{sd}},z^{\mathrm{sd}}}=z^{\mathrm{sd}}$.

\paragraph{RLT identities.}
Same pattern as~\S\ref{sec:els}: one per equality row.  Squaring the
state transition~\eqref{eq:hp-state} produces a 4-variable identity;
squaring each complement produces a 2-variable identity.

\subsection{Non-flexible demand (\texorpdfstring{$\textsc{d}_k$}{D\_k})}
\label{sec:demand}

\paragraph{Original constraints.}
\begin{align}
d^{\mathrm{nfl},k}_{i,t}&=\bar d^{k}_{i,t},
  &&\text{(demand fixing)}\label{eq:demand-fix}\\
0\le i^{\mathrm{gri},k}_{i,t}&\le\overline{I}^{k},
  &&\text{(import bound)}\label{eq:demand-import}
\end{align}

\paragraph{Equality form.}
Demand fixing is already an equality.  Import bound:
$i^{\mathrm{gri},k}_{i,t}+\sigma^{\mathrm{imp},k}_{i,t}
=\overline{I}^{k}$.

\paragraph{Binaries.} None.

\paragraph{Objective.}
$w_i$ entry: $+u^{k}_{i,t}$ on $d^{\mathrm{nfl},k}_{i,t}$ (consumer
utility, already positive in profit convention);
$-\pi^{k,\mathrm{imp}}_t$ on $i^{\mathrm{gri},k}_{i,t}$ (import cost).

\subsection{Grid export bounds}
\label{sec:grid-export}

Any module that includes a grid-export variable $e^{\mathrm{gri},k}_{i,t}$
also has an upper bound
$e^{\mathrm{gri},k}_{i,t}\le\overline{E}^{k}_i$,
which in equality form becomes
$e^{\mathrm{gri},k}_{i,t}+\sigma^{\mathrm{exp},k}_{i,t}
=\overline{E}^{k}_i$.  This produces one additional RLT identity per
export-period pair.

%======================================================================
\section{The CPP reformulation}
\label{sec:cpp}
%======================================================================

\subsection{Step 1: equality form in the $x$-space}

For any player $i$ with technology set $\Tech_i$, assemble $A_ix_i=b_i$
by stacking the rows from each module in Sections~\ref{sec:balance}--\ref{sec:grid-export}.
The slacked vector $x_i\ge0$ includes all original variables, all inequality
slacks, and all binary complement slacks.  Write
\begin{equation}
\feaseq_i:=\{x_i\ge0:A_ix_i=b_i\}.
\label{eq:feaseq}
\end{equation}

\subsection{Step 2: lifting to the $X$-space}
\label{sec:lifting}

Introduce $X_{ii}\approx x_ix_i^{\top}$.  The
\emph{RLT~identities} are the diagonal of the squared linear system:
\begin{equation}
\feasbar_i:=\bigl\{X_{ii}:\inner{a_{i,j}a_{i,j}^{\top}}{X_{ii}}
  =b_{i,j}^{2}\;\;\forall j\bigr\}
\quad\Longleftrightarrow\quad
\diag\!\bigl(A_iX_{ii}A_i^{\top}\bigr)=b_i\circ b_i,
\label{eq:rlt}
\end{equation}
where $\circ$ denotes the Hadamard (entrywise) product.  For each
binary index $m\in B_i$, the identity $z_{i,m}-z_{i,m}^{2}=0$ lifts to
\begin{equation}
(X_{ii})_{m,m}=z_{i,m}.
\label{eq:bindiag}
\end{equation}

\subsection{Step 3: per-player lifted set}

\begin{equation}
\liftset_i:=\bigl\{(x_i,X_{ii}):
  x_i\in\feaseq_i,\;X_{ii}\in\feasbar_i,\;
  (X_{ii})_{m,m}=z_{i,m}\;\forall m\in B_i\bigr\}.
\label{eq:Fi}
\end{equation}
This is purely affine (linear equalities plus the diagonal binary
identities); it carries no conic constraint.  For players without
binaries ($B_i=\varnothing$, e.g.\ a pure demand player), $\liftset_i$
reduces to the linear constraints plus RLT identities alone.

\subsection{Step 4: the CPP}

Assemble the joint matrix from the per-player blocks:
\[
Y_S=\begin{pmatrix}1 & x_S^{\top}\\ x_S & X_S\end{pmatrix},
\qquad X_S=(X_{ij})_{i,j\in S},
\]
with $x_S=(x_i)_{i\in S}$ and off-diagonal blocks $X_{ij}\approx
x_ix_j^{\top}$.  The chain:
\begin{align}
\vmip(S)
&=\max\Bigl\{\textstyle\sum_{i\in S}w_i^{\top}x_i\;:\;
  (x_i,X_{ii})\in\liftset_i\;\forall i,\;
  \textstyle\sum_i\Acom_i\xcom_i=0,\;
  X_S=x_Sx_S^{\top}\Bigr\}
\label{eq:chain-rank1}\\
&\le\max\Bigl\{\textstyle\sum_{i\in S}w_i^{\top}x_i\;:\;
  (x_i,X_{ii})\in\liftset_i\;\forall i,\;
  \textstyle\sum_i\Acom_i\xcom_i=0,\;
  \Tr\!\bigl((\Acom)^{\top}\Acom X_S\bigr)=0,\;
  Y_S\in\CP\Bigr\}
\label{eq:chain-cpp}\\
&=\vmip(S).
\label{eq:chain-eq}
\end{align}

Line~\eqref{eq:chain-rank1} re-expresses~\eqref{eq:vmip} with the
rank-one constraint $X_S=x_Sx_S^{\top}$.
Line~\eqref{eq:chain-cpp} relaxes rank-one to joint completely positive
membership $Y_S\in\CP$ and adjoins the squared coupling
$\Tr((\Acom)^{\top}\Acom X_S)=0$.
Line~\eqref{eq:chain-eq} is Burer's theorem (exactness, verified below).

The CPP is therefore:
\begin{equation}
\boxed{
\vmip(S)=\max_{Y_S}\;\inner{W_S}{Y_S}\quad
\text{s.t.}\quad
\begin{aligned}
&(x_i,X_{ii})\in\liftset_i&&\forall i\in S,\\
&\textstyle\sum_{i\in S}\Acom_i\xcom_i=0,\\
&\Tr\!\bigl((\Acom)^{\top}\Acom X_S\bigr)=0,\\
&Y_S\in\CP.
\end{aligned}}
\label{eq:cpp}
\end{equation}

\paragraph{Objective matrix.}
Since the objective $\sum_iw_i^{\top}x_i$ is linear,
\begin{equation}
W_i=\tfrac{1}{2}\begin{pmatrix}0&w_i^{\top}\\w_i&0\end{pmatrix},
\qquad\inner{W_i}{Y_i}=w_i^{\top}x_i,
\qquad W_S=\diag(W_i)_{i\in S}.
\label{eq:Wmat}
\end{equation}
No off-diagonal block in $W_S$: the objective has no inter-player product.

\subsection{Exactness (Burer's theorem)}

\begin{proposition}
\label{prop:exact}
The CPP~\eqref{eq:cpp} is an exact reformulation of~\eqref{eq:vmip}:
the equality~\eqref{eq:chain-eq} holds.
\end{proposition}

\begin{proof}
We verify the key assumption of Burer's theorem (Bomze--Gabl,
Thm.~4), module by module.

\emph{Condition~(a):} The only quadratic constraints are the binary
identities $z_{i,m}-z_{i,m}^2=0$.  On the slacked feasible set
$\feaseq_i$, each binary satisfies $z_{i,m}\in[0,1]$ (from
$z+\sigma=1$, $z,\sigma\ge0$), so
$z_{i,m}-z_{i,m}^2\ge0$~\checkmark.

\emph{Condition~(b):} The recession cone $L_{\infty}=\{x\ge0:Ax=0\}$.
Each technology module contributes finite upper bounds on every
variable: renewable availability, storage capacity, grid import/export
caps, electrolyzer capacity, heat pump capacity, and fixed demand
profiles.  In equality form these bounds become
$x_j+\sigma_j=\bar x_j$ with $\bar x_j>0$, so any $d\in L_{\infty}$
satisfies $d_j+d_{\sigma_j}=0$ with $d_j,d_{\sigma_j}\ge0$, forcing
$d_j=0$.  Hence $L_{\infty}=\{0\}$ and condition~(b) is
vacuous~\checkmark.

For players without binaries ($B_i=\varnothing$), there are no
quadratic constraints, so condition~(a) is vacuous;
condition~(b) still holds by boundedness.
\end{proof}

%======================================================================
\section{The Lagrangian dual (COP)}
\label{sec:dual}
%======================================================================

\subsection{Dualization}

Starting from the CPP~\eqref{eq:cpp}, dualize two constraint families:
\begin{align*}
\text{dualize:}\quad
&Y_S\in\CP &&\rightsquigarrow\; M\in\COP,\\
&\Tr\!\bigl((\Acom)^{\top}\Acom X_S\bigr)=0
  &&\rightsquigarrow\; \lambda\in\R;\\
\text{keep:}\quad
&(x_i,X_{ii})\in\liftset_i\;\forall i,
\qquad\textstyle\sum_i\Acom_i\xcom_i=0.
\end{align*}
The Lagrangian characteristic function is
\begin{equation}
\vlr_S(M,\lambda)=\max_{\substack{(x_i,X_{ii})\in\liftset_i\;\forall i\in S\\
\sum_{i\in S}\Acom_i\xcom_i=0}}\;
\inner{W_S+M}{Y_S}-\lambda\,\Tr\!\bigl((\Acom)^{\top}\Acom X_S\bigr).
\label{eq:lagrangian}
\end{equation}
The dual is
$\min_{M\in\COP,\,\lambda}\vlr_S(M,\lambda)$.

\subsection{Cross-block vanishing}

The off-diagonal blocks $X_{ij}$ ($i\neq j$) are \emph{free}: after
dualization, no constraint reaches them ($\liftset_i$ touches only
$X_{ii}$; the linear coupling touches only $\xcom_i$).  A free variable
with nonzero linear coefficient sends the maximization to $+\infty$;
finiteness forces:

\begin{lemma}[Finiteness conditions]
\label{lem:vanish}
$\vlr_S(M,\lambda)<\infty$ only if, for every $i\neq j$ in $S$:
\begin{align}
(M_{ij})_{\mathrm{cc}}&=\lambda\,(\Acom_i)^{\top}\Acom_j,
\label{eq:fin-cc}\\
(M_{ij})_{\mathrm{cr}}&=(M_{ij})_{\mathrm{rc}}
=(M_{ij})_{\mathrm{rr}}=0,
\label{eq:fin-rest}
\end{align}
where the subscripts $\mathrm{c}$,$\mathrm{r}$ refer to the
coupling/rest partition of the variable indices.
\end{lemma}

\paragraph{Carrier-wise sparsity.}
The product $(\Acom_i)^{\top}\Acom_j$ sums over the $R$ coupling rows.
A row for carrier $k$, period $t$ contributes a nonzero term only if
\emph{both} players $i$ and $j$ have coupling variables in carrier $k$
(i.e.\ both participate in the $k$-community balance).  Concretely:
\begin{itemize}[leftmargin=1.2em]
\item If $i$ and $j$ share \emph{no} carrier:
  $(\Acom_i)^{\top}\Acom_j=0$ and $M_{ij}=0$ entirely.
\item If $i$ and $j$ both participate in carrier $k$: the product has
  nonzero entries in the rows/columns corresponding to their
  $k$-coupling variables, summed over the $T$ periods.
\end{itemize}
This sparsity is technology-determined: two demand-only players in
different carriers have zero cross-block.

\subsection{Reduced LPG}

Under the finiteness conditions~\eqref{eq:fin-cc}--\eqref{eq:fin-rest},
the cross blocks contribute nothing to the value, and
\begin{equation}
\boxed{
\vlr_S(M,\lambda)=\max_{\substack{(x_i,X_{ii})\in\liftset_i\;\forall i\in S\\
\sum_i\Acom_i\xcom_i=0}}\;
\sum_{i\in S}\inner{\widetilde W_i}{Y_i},
\qquad
\widetilde W_i:=W_i+M_{ii}
  -\lambda\,(\Acom_i)^{\top}\Acom_i\,\Pi_{\mathrm{cc}},}
\label{eq:reduced}
\end{equation}
where $(\Acom_i)^{\top}\Acom_i$ sits on the coupling--coupling sub-block
via $\Pi_{\mathrm{cc}}$.  This is a \emph{classical linear production
game} (affine resource sets, linear objective, one homogeneous coupling)
with a nonempty core by Owen's theorem.

\paragraph{Owen allocation.}
The dual of the grand-coalition reduced LPG ($S=\N$) yields multipliers
$(\mu^{*},\{\nu_i^{*}\})$.  The allocation
\begin{equation}
\alpha^{\mathrm{Owen}}_i:=\inner{\nu_i^{*}}{g_i},\qquad i\in\N,
\label{eq:owen}
\end{equation}
lies in the core: $\sum_{i\in S}\alpha^{\mathrm{Owen}}_i
\ge\vlr_S(M,\lambda)$ for all $S$, with equality at $S=\N$.

%======================================================================
\section{Main theorem}
\label{sec:theorem}
%======================================================================

\begin{theorem}[Core of the mixed-integer community game]
\label{thm:main}
Assume:
\begin{enumerate}[label=\textup{(A\arabic*)}]
\item \textbf{Zero duality gap} (automatic from boundedness):
  $\inf_{M\in\COP,\lambda}\vlr_\N(M,\lambda)=\vmip(\N)$.
\item \textbf{Dual attainment} (instance-dependent hypothesis):
  the infimum is achieved at some $(M^{*},\lambda^{*})$.
\end{enumerate}
Then:
\begin{enumerate}[label=\textup{(\roman*)}]
\item \emph{Exact core} (conditional on~\textup{(A2)}):
  the Owen allocation $\chi:=\alpha^{\mathrm{Owen}}$ at
  $(M^{*},\lambda^{*})$ satisfies
  \[
  \sum_{i\in S}\chi_i\ge\vmip(S)\;\;\forall S\subseteq\N,
  \qquad\sum_{i\in\N}\chi_i=\vmip(\N).
  \]
\item \emph{$\varepsilon$-core} (unconditional): for every
  $\varepsilon>0$, the Owen allocation at an $\varepsilon$-optimal
  multiplier satisfies stability and
  $\vmip(\N)\le\sum_{i\in\N}\chi^{\varepsilon}_i
  \le\vmip(\N)+\varepsilon$.
\end{enumerate}
\end{theorem}

The proof is the composition of Corollary~\ref{cor:upper} (coalition-wise
upper bound from weak duality) with the Owen core of the reduced LPG
(Section~\ref{sec:dual}).

\begin{corollary}[Coalition-wise upper bound]
\label{cor:upper}
For every $S\subseteq\N$ and every $(M,\lambda)$ with $M\in\COP$
satisfying~\eqref{eq:fin-cc}--\eqref{eq:fin-rest}:
$\vlr_S(M,\lambda)\ge\vmip(S)$.
\end{corollary}

%======================================================================
\section{Companion identity: per-player lift = convex hull pricing}
\label{sec:chp}
%======================================================================

If one lifts only the per-player feasible sets (imposing $Y_i\in\CP_i$
and $\liftset_i$ block by block) while keeping the linear coupling and
dropping the squared coupling and joint cone:
\begin{equation}
\mathrm{(P)}:\quad
\max_{\{Y_i\}}\;\sum_{i\in\N}\inner{W_i}{Y_i}
\quad\text{s.t.}\quad
\sum_{i\in\N}\Acom_i\xcom_i=0,\quad
(x_i,X_{ii})\in\liftset_i,\;Y_i\in\CP_i\;\;\forall i,
\label{eq:chp-lift}
\end{equation}
then (P) has the same optimal value and the same community-balance
multipliers as convex hull pricing:
\begin{equation}
\vchp(\N)=\max_{\{x_i\}}\;\sum_{i\in\N}w_i^{\top}x_i
\quad\text{s.t.}\quad
\sum_{i\in\N}\Acom_i\xcom_i=0,\;\;
x_i\in\conv(\feas_i)\;\;\forall i.
\label{eq:chp}
\end{equation}
This follows from the per-player projection
$\mathrm{proj}_{x_i}\{(x_i,X_{ii})\in\liftset_i,Y_i\in\CP_i\}
=\conv(\feas_i)$ (Burer at the per-player level).

%======================================================================
\section{Deferred extensions (future work)}
\label{sec:deferred}
%======================================================================

Two additional coupling families from the preprint are deferred:

\paragraph{Reserve adequacy.}
$r^{\mathrm{sym}}-\sum_{i\in S}\sum_k r^{k,+}_{i,t}\le0$ (per period,
homogeneous).  This introduces a common variable
$x_0=(r^{\mathrm{sym}})$ controlled by whichever coalition operates,
not owned by any player.  The coupling becomes
$\sum_{i\in S}\Acom_i\xcom_i+A_0x_0=0$, and the storage headroom
variables $r^{k,\pm}_{i,t}$ join the coupling coordinates.  Cross-block
vanishing extends to prosumer--common blocks $X_{i0}$ by the same
mechanism (note~revised, Lemma~4).

\paragraph{Peak penalty.}
$\sum_{i\in S}\tfrac{1}{\Delta T}(i^{E,\mathrm{mkt}}_{i,t}
-e^{E,\mathrm{mkt}}_{i,t})-p\le0$ (per period, homogeneous).  Common
variable $p$ (peak import power); grid-exchange variables become
coupling coordinates.

Both constraints are homogeneous (RHS$\,=0$), so the LPG structure is
preserved.  They are excluded here to keep the initial implementation
focused on the community balance coupling alone.

%======================================================================
\section{Computational notes}
\label{sec:comp}
%======================================================================

\paragraph{Dimension of the lifted matrix.}
$Y_S$ has dimension $1+\sum_{i\in S}n_i^{\mathrm{eq}}$, where
$n_i^{\mathrm{eq}}=n_i+(\text{\# inequality slacks in }
\feas_i)+|B_i|$ is the slacked dimension of player~$i$.  The exact
count is technology-determined:
\begin{center}\small
\begin{tabular}{@{}l c c c@{}}
\toprule
Module & Continuous per $t$ & Binaries per $t$ & Ineq.\ slacks per $t$ (approx.)\\
\midrule
\textsc{res} & 2 & 0 & 1 \\
$\textsc{ess}_k$ & 3 & 0 & 3 \\
\textsc{els} ($N_s=6$) & $3+N_s=9$ & $5+N_s=11$ & $\sim 2N_s+8=20$ \\
\textsc{hp} & 4 & 3 & $\sim 5$ \\
$\textsc{d}_k$ / $\textsc{fl}_k$ & 2 & 0 & 1 \\
\bottomrule
\end{tabular}
\end{center}
For the standard 6-player community ($T=24$):
$\sum_in_i^{\mathrm{eq}}\approx2000$--$2500$ (the precise count depends
on wrap-around and downtime constraints), giving $Y_S$ of dimension
$\sim\!2500$.

\paragraph{Technology modularity.}
Adding a player with technology set $\Tech_i$ increases the dimension by
$n_i^{\mathrm{eq}}$ (determined by $\Tech_i$ alone) and adds one block
to $Y_S$.  The coupling matrix $\Acom$ gains columns for the new
player's coupling variables.  No existing block is modified.  Scaling
from 6 to 30 players is therefore mechanical: assemble the per-player
modules and stack.

\paragraph{Cutting-plane algorithm.}
The COP dual is solved by outer approximation.
Section~\ref{sec:algorithm} presents the baseline algorithm
(Guo~et~al., 2024) and the enhanced set-copositive version
(Gabl~\&~Anstreicher, 2025) that we adopt.

%======================================================================
\section{Algorithm design: cutting-plane outer approximation}
\label{sec:algorithm}
%======================================================================

\subsection{The CPP in CPOPT standard form}
\label{sec:cpopt}

Both algorithms operate on the same standard-form program.  Define the
joint homogenised variable
\begin{equation}
y:=\begin{pmatrix}y_0\\x_S\end{pmatrix}\in\R^{1+n},
\qquad y_0=1,\quad n=\textstyle\sum_{i\in S}n_i^{\mathrm{eq}},
\label{eq:homog-y}
\end{equation}
and the lifted matrix $Y=yy^{\top}\in\R^{(1+n)\times(1+n)}$, relaxed
to $Y\in\mathrm{CP}^{1+n}$.  Every affine constraint of the
CPP~\eqref{eq:cpp} becomes a linear constraint on~$Y$:

\begin{enumerate}[label=(\alph*),leftmargin=2em]
\item \emph{Per-player linear constraints}
  $A_ix_i=b_i$ (from $\feaseq_i$): homogenise as
  $\bar A_iy=0$ with $\bar A_i=[-b_i\;|\;0\;\cdots\;A_i\;\cdots\;0]$,
  giving first-column constraints $\bar A_iYe_0=0$.

\item \emph{Per-player RLT identities}
  $\inner{a_{i,j}a_{i,j}^{\top}}{X_{ii}}=b_{i,j}^{2}$: the squared
  form of~(a), equivalently
  $\inner{\bar a_{i,j}\bar a_{i,j}^{\top}}{Y}=0$.

\item \emph{Binary diagonal}
  $(X_{ii})_{m,m}=z_{i,m}$ ($m\in B_i$): $Y_{m,m}-Y_{0,m}=0$.

\item \emph{Linear coupling}
  $\sum_i\Acom_i\xcom_i=0$: homogenise as
  $\bar A^{\mathrm{cpl}}y=0$.

\item \emph{Squared coupling}
  $\Tr((\Acom)^{\top}\Acom X_S)=0$: the squared form of (d),
  $\inner{\bar a^{\mathrm{cpl}}_r(\bar a^{\mathrm{cpl}}_r)^{\top}}{Y}=0$.

\item \emph{Normalisation} $Y_{00}=1$.
\end{enumerate}

Collecting as $B_l\bullet Y=b_l$, $l=1,\dots,m$:
\begin{equation}
\boxed{
\mathrm{CPOPT}:\quad
\min_{Y}\;C\bullet Y
\quad\text{s.t.}\quad
B_l\bullet Y=b_l,\;l=1,\dots,m,\qquad
Y\in\mathrm{CP}^{1+n},}
\label{eq:cpopt}
\end{equation}
with $C=-W_S$ (sign flip: CPOPT minimises, our game maximises profit).
Its conic dual is:
\begin{equation}
\mathrm{COPOPT}:\quad
\max_{v\in\R^m}\;b^{\top}v
\quad\text{s.t.}\quad
S:=C-\textstyle\sum_{l=1}^{m}v_lB_l\in\mathrm{COP}_{1+n}.
\label{eq:copopt}
\end{equation}
Both algorithms solve COPOPT by outer-approximating the copositive
constraint on~$S$.  They differ in the separation oracle.

%----------------------------------------------------------------------
\subsection{Baseline: Guo~et~al.\ cutting-plane (MILPCOP)}
\label{sec:guo}
%----------------------------------------------------------------------

\paragraph{Master problem.}
Remove $S\in\mathrm{COP}_{1+n}$ from~\eqref{eq:copopt} and replace it
by an outer approximation $\mathcal{C}\supseteq\mathrm{COP}_{1+n}$,
initialised with:
\begin{enumerate}[leftmargin=1.6em]
\item Diagonal nonnegativity: $S_{jj}\ge0\;\forall j$.
\item SOCP (DNN) strengthening: decompose $S=V+N$, $N\ge0$, with
  $2\times2$ PSD relaxation $V_{jj}\ge0$,
  $V_{jj}V_{kk}\ge V_{jk}^{2}$ for all $j<k$ (rotated SOC).
\end{enumerate}
Solve the master as an SOCP (or LP if the DNN part is omitted).

\paragraph{Separation oracle: MILPCOP (Anstreicher, 2021).}
Given candidate $\bar S$, test $\bar S\in\mathrm{COP}_{1+n}$:
\begin{equation}
\mathrm{MILPCOP}:\quad
\max_{\gamma,u,z}\;\gamma
\quad\text{s.t.}\quad
\bar S\,u\le-\gamma\,e+\nu\circ(e-z),\;
0\le\gamma,\;
0\le u\le z,\;
z\in\{0,1\}^{1+n},\;
e^{\top}z\ge\theta,
\label{eq:milpcop-guo}
\end{equation}
with big-$M$ coefficients
$\nu_j=1+\sum_{k\neq j}\max(0,\bar S_{jk})$ and $\theta=2$ (since
$\diag(S)\ge0$ is enforced).

\emph{Certificate:} $\gamma^{*}=0\Leftrightarrow\bar S\in
\mathrm{COP}_{1+n}$.  If $\gamma^{*}>0$, the witness $\bar u\ge0$
satisfies $\bar u^{\top}\bar S\bar u<0$; add cut
$\bar u^{\top}S\bar u\ge0$ to the master.

\paragraph{Iteration.}
Repeat master$\to$separation$\to$cut until $\gamma^{*}=0$ or
time/iteration limit.  The master objective is non-increasing.

\paragraph{Limitation.}
The witness $\bar u$ in~\eqref{eq:milpcop-guo} is constrained only by
$u\ge0$ --- it need not satisfy the energy balance, SOC transitions, or
any other structure of the Burer lift.  In practice this produces
\emph{wasted cuts}: valid inequalities that do not improve the bound,
because their witness vector has no relation to the primal feasible
region.  Guo~et~al.\ report convergence in $\sim\!60$--$1600$
iterations for $n_c\le57$ but $>3000$ iterations (often without
convergence) for $n_c\ge113$.

%----------------------------------------------------------------------
\subsection{Enhanced: Gabl~\&~Anstreicher set-copositive cutting-plane
(MILPCOP($K$))}
\label{sec:ga}
%----------------------------------------------------------------------

The key modification: replace $\mathrm{COP}_{1+n}$ with the
\emph{set-copositive cone} $\mathrm{COP}(K)$, where
$K=\{y\ge0:\mathcal{A}y=0\}$ encodes the first-order equality
structure of the Burer lift.  This restricts the separation oracle to
witnesses that respect the problem structure.

\subsubsection{Constraint partition: what goes into $K$ vs.\ CPOPT}
\label{sec:partition}

The CPOPT constraints fall into two classes:
\begin{enumerate}[label=\textbf{(\Alph*)},leftmargin=2.4em]
\item \textbf{First-order (linear in $y$):}  $\bar a_l^{\top}y=0$.
  These are homogenised versions of the original linear equalities
  $A_ix_i=b_i$ (per-player) and
  $\sum_i\Acom_i\xcom_i=0$ (coupling), plus normalisation $y_0=1$.
  \emph{These define the ground cone~$K$.}
\item \textbf{Lifted (quadratic in $y$, linear in $Y$):}
  RLT identities $\inner{\bar a\bar a^{\top}}{Y}=0$, binary diagonals
  $Y_{mm}-Y_{0m}=0$, squared coupling
  $\inner{\bar a^{\mathrm{cpl}}_r(\bar a^{\mathrm{cpl}}_r)^{\top}}{Y}=0$.
  These are CPOPT constraints but do \emph{not} enter~$K$.
\end{enumerate}

\begin{remark}[Why lifted constraints stay out of $K$]
\label{rem:lifted-not-K}
For rank-one $Y=yy^{\top}$ with $y\in K$, every Class~B constraint
is automatically satisfied: RLT from $\bar a^{\top}y=0$, binary
diagonal from $y_m\in\{0,1\}$ (encoded by $y_m+\sigma_m=1$ in~$K$
plus the rank-one property), squared coupling from the linear coupling.
Class~B constraints are redundant on rank-one solutions; they tighten
the relaxation only when $Y\neq yy^{\top}$.  Adding them to~$K$ would
shrink~$K$, enlarge $\mathrm{COP}(K)$, and \emph{weaken} the outer
approximation.
\end{remark}

\paragraph{The CPOPT constraint list.}
\begin{center}\small
\begin{tabular}{@{}l l l@{}}
\toprule
Constraint & Class & Count\\
\midrule
Per-player linear $\bar A_iy=0$ & A ($\to K$) & $\sum_i p_i$\\
Coupling $\bar A^{\mathrm{cpl}}y=0$ & A ($\to K$) & $R=3T$\\
Normalisation $y_0=1$ & A ($\to K$) & $1$\\[2pt]
Per-player RLT $\inner{\bar a_{i,j}\bar a_{i,j}^{\top}}{Y}=0$
  & B (CPOPT only) & $\sum_i p_i$\\
Binary diagonal $Y_{mm}-Y_{0m}=0$ & B (CPOPT only) & $\sum_i|B_i|$\\
Squared coupling
  $\inner{\bar a^{\mathrm{cpl}}_r(\bar a^{\mathrm{cpl}}_r)^{\top}}{Y}=0$
  & B (CPOPT only) & $R=3T$\\
\bottomrule
\end{tabular}
\end{center}
Total: $m=2(\sum_ip_i+R)+\sum_i|B_i|+1$.

\subsubsection{The ground cone $K$}
\label{sec:ground-cone}

Collect the Class~A rows into
$\mathcal{A}\in\R^{p\times(1+n)}$:
\begin{equation}
K:=\bigl\{y\in\R^{1+n}_+:\mathcal{A}\,y=0\bigr\}.
\label{eq:ground-cone}
\end{equation}
$\mathcal{A}$ stacks:
\begin{enumerate}[leftmargin=1.6em]
\item \emph{Per-player first-order equalities} $\bar A_i$:
  energy balance, SOC transitions, device coupling, state selection,
  power-consumption coupling, demand fixing, bound slacks
  ($x+\sigma=\bar x$), binary complements ($z+\sigma^{z}=1$)
  --- \textbf{block-diagonal};
\item \emph{Coupling} $\bar A^{\mathrm{cpl}}$: the $R=3T$
  community-balance rows --- \textbf{only inter-block rows};
\item \emph{Normalisation} $e_0^{\top}y=1$.
\end{enumerate}
Row count: $p=\sum_i p_i+R+1$.  Extremely sparse (item~1
block-diagonal, item~2 at most $|\N|$ nonzeros per row).

\subsubsection{The set-copositive COPOPT}

Replace $\mathrm{COP}_{1+n}$ in~\eqref{eq:copopt} by
$\mathrm{COP}(K)=\{M\in\mathcal{S}^{1+n}:y^{\top}My\ge0\;\forall y\in K\}$:
\begin{equation}
\mathrm{COPOPT}(K):\quad
\max_{v\in\R^m}\;b^{\top}v
\quad\text{s.t.}\quad
S:=C-\textstyle\sum_{l=1}^{m}v_lB_l\in\mathrm{COP}(K).
\label{eq:copopt-K}
\end{equation}
Since $\mathrm{COP}_{1+n}\subset\mathrm{COP}(K)$, the feasible set
is \emph{larger}, but Gabl~\&~Anstreicher Lemma~1 ensures
$\mathrm{val}(\mathrm{COPOPT}(K))=\mathrm{val}(\mathrm{COPOPT})$
whenever the CPOPT has a rank-one optimum (guaranteed by Burer's
theorem and boundedness).

\paragraph{Strong duality.}
By Cifuentes--Dey--Xu (2025, Thm.~3), boundedness of $\feas_i$ implies
$\mathrm{val}(\mathrm{CPOPT})=\mathrm{val}(\mathrm{COPOPT})=\vmip(S)$.

\paragraph{Dual variable interpretation.}
$v=(v_1,\dots,v_m)$ maps to:
$v_l$ for per-player constraint $\leftrightarrow\nu_{i,j}$;
$v_l$ for coupling $\leftrightarrow\mu_r$ (community price);
$v_l$ for RLT $\leftrightarrow$ (RLT dual);
$v_l$ for squared coupling $\leftrightarrow\lambda_r$;
$v_l$ for binary diagonal $\leftrightarrow\delta_m$.

\subsubsection{MILPCOP($K$) separation oracle}
\label{sec:separation}

Given candidate $\bar S$ from the master, test
$\bar S\in\mathrm{COP}(K)$:
\begin{equation}
\boxed{
\mathrm{MILPCOP}(K):\quad
\begin{aligned}
\max_{\gamma,u,z,\lambda}\quad&\gamma\\
\text{s.t.}\quad
&\bar S\,u+\mathcal{A}^{\top}\lambda
  \le-\gamma\,e+\nu\circ(e-z),\\
&\mathcal{A}\,u=0,\\
&0\le\gamma\le\alpha\,e^{\top}u,\\
&0\le u\le z,\\
&z\in\{0,1\}^{1+n},
\end{aligned}}
\label{eq:milpcop-K}
\end{equation}
where $\lambda\in\R^{p}$ are dual multipliers for
$\mathcal{A}u=0$, $\nu_j=1+\sum_{k\neq j}\max(0,\bar S_{jk})$,
$\alpha>0$.

\emph{Certificate:} $\gamma^{*}=0\Leftrightarrow
\bar S\in\mathrm{COP}(K)$ (Gabl~\&~Anstreicher, Lemma~2).
If $\gamma^{*}>0$, the witness $\bar u\in K$ satisfies
$\bar u^{\top}\bar S\bar u<0$; add cut $\bar u^{\top}S\bar u\ge0$.

\subsubsection{What changes from Guo~et~al.\ to Gabl~\&~Anstreicher}

\begin{center}\small
\begin{tabular}{@{}p{0.22\textwidth} p{0.34\textwidth} p{0.34\textwidth}@{}}
\toprule
& Guo et al.\ (\S\ref{sec:guo}) & Gabl \& Anstreicher (\S\ref{sec:ga})\\
\midrule
Cone & $\mathrm{COP}_{1+n}$ & $\mathrm{COP}(K)$, $K=\{y\ge0:\mathcal{A}y=0\}$\\
Separation MILP & \eqref{eq:milpcop-guo}: $u\ge0$ only
  & \eqref{eq:milpcop-K}: $u\ge0$ \emph{and} $\mathcal{A}u=0$\\
Extra variables & --- & $\lambda\in\R^{p}$ (dual of $\mathcal{A}u=0$)\\
Cut witness & Any $u\ge0$ with $u^{\top}\bar Su<0$
  & Only $u\in K$ with $u^{\top}\bar Su<0$\\
Cut relevance & May violate energy bal., SOC, etc.\
  & Respects all first-order structure\\
Iteration count & $\sim\!1000$+ (Guo, Table~6--7)
  & Fewer (structure-aware cuts)\\
Optimality & $S^{*}\in\mathrm{COP}_{1+n}$
  & $S^{*}\in\mathrm{COP}(K)$ (same opt.\ value by Lemma~1)\\
\bottomrule
\end{tabular}
\end{center}

\begin{remark}[Relationship to note~revised]
In the note~revised framework
(Section~\ref{sec:dual}), the squared coupling is dualised with
multiplier~$\lambda$, and cross-block vanishing gives the reduced LPG.
In the Gabl~\&~Anstreicher framework, the squared coupling is a Class~B
CPOPT constraint, and cross-block vanishing is a structural property of
the optimal $S^{*}$ but is not needed algorithmically.  The two are
equivalent by strong duality; we implement the latter for the superior
separation oracle.
\end{remark}

%----------------------------------------------------------------------
\subsection{Master problem initialisation}
\label{sec:master-init}
%----------------------------------------------------------------------

The master replaces $S\in\mathrm{COP}(K)$ by an outer approximation
$\mathcal{C}\supseteq\mathrm{COP}(K)$:
\begin{enumerate}[leftmargin=1.6em]
\item \emph{Diagonal nonnegativity}: $S_{jj}\ge0\;\forall j$.
\item \emph{$2\times2$ SDD constraints} (SOCP): for all $j<k$,
  \begin{equation}
  \bar S_{jk}^{2}\le S_{jj}\,S_{kk},\quad \bar S_{jk}\ge S_{jk}.
  \label{eq:sdd}
  \end{equation}
\item \emph{DNN decomposition} (optional): $V+N=S$, $N\ge0$,
  $V_{jj}\ge0$, $V_{jj}V_{kk}\ge V_{jk}^{2}$.
\end{enumerate}

%----------------------------------------------------------------------
\subsection{The full algorithm (pseudocode)}
\label{sec:full-algo}
%----------------------------------------------------------------------

\begin{equation}
\boxed{
\begin{aligned}
&\textbf{Algorithm: Set-Copositive Outer Approximation for the Community COP}\\
&\text{---}\\
&\textbf{Input:}\; (C,\{B_l,b_l\},\mathcal{A});\;
  \varepsilon_{\gamma}>0;\;K_{\max},T_{\max}.\\
&\textbf{Output:}\; v^{*},\;S^{*},\;\chi.\\[4pt]
&1.\;\textbf{Initialise.}\;
  \mathcal{C}\leftarrow\{\diag(S)\ge0\}\cap\eqref{eq:sdd}.\\
&2.\;\textbf{For } k=0,1,2,\dots:\\
&\quad 2a.\;\text{Solve master }
  \mathrm{COPOPT}(K)|_{S\in\mathcal{C}}
  \;\to\;(v^{k},S^{k}).\\
&\quad 2b.\;\text{Compute } \nu_j\leftarrow
  1+\textstyle\sum_{l\neq j}\max(0,S^{k}_{jl}).\\
&\quad 2c.\;\text{Solve }\mathrm{MILPCOP}(K)
  \text{ for }\bar S=S^{k}
  \;\to\;(\gamma^{k},u^{k},z^{k},\lambda^{k}).\\
&\quad 2d.\;\textbf{If } \gamma^{k}\le\varepsilon_{\gamma}:\;
  v^{*}\!\leftarrow v^{k};\;S^{*}\!\leftarrow S^{k};\;
  \textbf{go to 3.}\\
&\quad 2e.\;\text{Add cut }(u^{k})^{\top}S\,u^{k}\ge0
  \text{ to }\mathcal{C}.\\
&\quad 2f.\;\textbf{If } k\ge K_{\max}\text{ or time}\ge T_{\max}:\;
  \text{return }v^{k}\text{ as }\varepsilon\text{-optimal.}\\
&3.\;\textbf{Extract allocation.}\\
&\quad \mu^{*}_r\leftarrow v^{*}_l\text{ for coupling row }r
  \quad\text{(community prices)}.\\
&\quad \nu^{*}_{i,j}\leftarrow v^{*}_l\text{ for per-player row }j
  \text{ of player }i.\\
&\quad \chi_i\leftarrow\textstyle\sum_j\nu^{*}_{i,j}\,b_{i,j}
  \quad\forall i\in\N
  \quad\text{(Owen allocation)}.
\end{aligned}}
\label{eq:algorithm}
\end{equation}

\subsection{Allocation extraction}
\label{sec:allocation}

At termination, the Owen allocation is:
\begin{equation}
\chi_i=\inner{\nu_i^{*}}{g_i}=
\textstyle\sum_j\nu^{*}_{i,j}\,b_{i,j},
\qquad i\in\N,
\label{eq:allocation-extract}
\end{equation}
where $\nu_i^{*}$ are the duals of the per-player constraints
(Class~A, item~1 of $\mathcal{A}$) and $g_i=(b_{i,j})$ the
right-hand sides.  The community prices $\mu^{*}$ (duals of the
coupling rows) give the internal energy prices.

If the algorithm terminates at an $\varepsilon$-optimal iterate,
$\chi$ is an $\varepsilon$-core allocation
(Theorem~\ref{thm:main}(ii)).

\subsection{Exploiting sparsity}
\label{sec:sparsity}

MILPCOP($K$)~\eqref{eq:milpcop-K} has $1+n$ binary variables $z$, but
$\mathcal{A}u=0$ with block-diagonal $\mathcal{A}$ forces strong
structural restrictions:

\begin{itemize}[leftmargin=1.2em]
\item \emph{Per-player nullspace.}  Within block~$i$,
  $\bar A_iu_i=0$ with $u_i\ge0$ restricts $u_i$ to the nullspace of
  $\bar A_i$ intersected with the nonneg orthant.  For a demand-only
  player ($\textsc{d}_k$), $\bar A_i$ nearly pins all variables, so
  the nullspace dimension is very small.
\item \emph{Coupling links are few.}  Only $3T$ coupling rows span
  blocks (vs.\ $n\sim2500$ columns).  The effective degrees of freedom
  are the per-player nullspace dimensions plus the coupling slack.
\item \emph{Preprocessing.}  Before solving MILPCOP($K$), compute the
  nullspace dimension of $\bar A_i$ for each block; fix $z_j=0$ for
  indices $j$ that cannot participate in any nullspace vector.
\end{itemize}

\subsection{Progressive scaling strategy}
\label{sec:scaling}

\begin{enumerate}[leftmargin=1.6em]
\item \emph{Phase~1 (validation):} 2--3 players, $T=4$--$8$.
  $n_c\le100$.  Verify against known $\vmip$ values.
\item \emph{Phase~2 (6-player):} Full model, $T=24$.
  $n_c\sim2500$.  Exploit sparsity; accept $\varepsilon$-core.
\item \emph{Phase~3 (scale-up):} 15--30 players, $T=24$.
  Block-wise preprocessing; parallel separation.
\end{enumerate}

\end{document}
